\documentclass{article}
\usepackage{polski}
\usepackage[utf8]{inputenc}
\usepackage{listings}
\usepackage{graphicx}
\usepackage{amsmath}  % dla matematycznych symboli

\title{AiSD Lista 3}
\author{Dominik Kowalczyk}

\begin{document}
	\maketitle
	
\section{Algorytm CUT ROD rozwiązujący problem podziału pręta}

\large{Chcemy znaleźć maksymalny zysk po jakim możemy sprzedać pręt długości n, mając możliwość dzielić go na kawałki. W celu rozwiązania tego problemu optymalizacyjnego może pomóc nam algorytm rozcinania pręta - CUT ROD, które przedstawimy implementację naiwną, ze spamiętywaniem oraz iteracyjną.}

\subsection{Naiwny}

\begin{lstlisting}[language=C++, caption={Kod funkcji rozcinania pręta w wersji naiwnej}, numbers=left, frame=single, xleftmargin=20pt]
float CUT_ROD_NAIWNY(float* p, int n) {
	if (n == 0) {
		return 0;
	}
	float q = inf;
	for (int i = 1; i <= n; ++i) {
		q = max(q, p[i - 1] + 
		CUT_ROD_NAIWNY(p, n - i));
	}
	return q;
}
\end{lstlisting}

\large {W naszym algorytmie funkcja $CUT ROD$ przyjmuje argumenty $p$, który jest wskaźnikiem przechowującym ceny różnych długości pręta oraz $n$ reprezentująca długość pręta.}

\large{Na początku wchodzimy do pętli $for$ sprawdzająca, czy nasz pręt nie jest długości $0$ (jeśli jest zwracamy $0$). Następnie tworzymy zmienną $q$ przechowującą maksymalny zysk i na start ustawiamy ją na \(-\infty\). Następnie przechodzimy do pętli $for$ przechodzącej po wszystkich możliwych podziałach pręta. Następnie obliczamy dla każdego podziału maksymalny zysk poprzez wybór za pomocą funkcji $max$ zysk aktualnie zapisany pod zmienną $q$ lub ceną odcinka długości $i$ lub pozostałej części pręta $n - i$ poprzez rekurencyjne wywołanie $CUT ROD NAIWNY$. Na sam koniec zwracamy $q$.}

\subsection{Ze spamiętywaniem}

\begin{lstlisting}[language=C++, caption={Kod funkcji rozcinania pręta w wersji ze spamiętywaniem}, numbers=left, frame=single, xleftmargin=20pt]
float MEMORIZED_CUT_ROD(float* p, float* r, 
int* s, int n) {
	if (r[n] >= 0) {
		return r[n];
	}
	float q;
	if (n == 0) {
		q = 0;
	} else {
		q = inf;
		
		for (int i = 1; i <= n; ++i) {
			
			float m = p[i - 1] +
			 MEMORIZED_CUT_ROD(p, 
			 r, s, n - i);
			 
			if (m > q) {
				q = m;
				s[n] = i;
			}
		}
	}
	r[n] = q;
	return q;
}
\end{lstlisting}

\large{W przypadku algorytmu ze spamiętywaniem mamy do czynienia z większą ilością argumentów. Oprócz $p$ i $n$ mamy wskaźnik (do tablicy) $r$ gdzie przechowywane jest maksymalny możliwy zysk dla danej długości pręta, jak również wskaźnik (do tablicy) $s$, który przechowuje informacje o optymalnych cięciach pręta.}

\large{Mamy warunki początkowe, które działają na tej samej zasadzie co w algorytmie naiwnym. Znacząca zmiana zachodzi w pętli $for$ gdzie tworzymy zmienną tymczasową $m$ wyliczając dla każdego cięcia $i$ zysk. Następnie mamy krótką pętle $if$, która przy warunku, gdy nasz zysk wyliczony i przechowywany w zmiennej tymczasowej $m$ jest większy od $q$ (zysku maksymalnego) to staje się nim, a nasza tablica optymalnego cięcia zapamiętuje wartość $i$. Ponownie zwracamy nasz zysk maksymalny.}

\subsection{iteracyjna}

\begin{lstlisting}[language=C++, caption={Kod funkcji rozcinania pręta w wersji ze spamiętywaniem}, numbers=left, frame=single, xleftmargin=20pt]
float EXT_BOT_UP_CUT_ROD(float* p, float* r,
 int* s, int n) {
	r[0] = 0;
	for (int j = 1; j <= n; ++j) {
		float q = inf;
		
		for (int i = 1; i <= j; ++i) {
			if (q < p[i - 1] + 
			 r[j - i]) {
			 	
				q = p[i - 1] + 
				 r[j - i];
				 
				s[j] = i;
			}
		}
		r[j] = q;
	}
	return r[n];
}
\end{lstlisting}

\large{W wersji iteracyjnej mamy takie same argumenty jak w przypadku ze spamiętywaniem: $p$, $r$, $s$, $n$. Początkowo ustawiamy maksymalny zysk na $0$ ($r[0] = 0)$. Wchodzimy do pętli $for$, ustawiamy ponownie $q$ na \(-\infty\). Wchodzimy do wewnętrznej pętli $for$, gdzie iterujemy przez wszystkie możliwe cięcia pręta.}

\begin{enumerate}
	\item Dla każdego cięcia obliczana jest suma:
	\[
	\text{zysk} = p[i - 1] + r[j - i],
	\]
	gdzie:
	\begin{itemize}
		\item $p[i - 1]$ - cena odciętego kawałka pręta o długości $i$ 
		\item $r[j - i]$ - maksymalny zysk dla pozostałej długości pręta $j - i$, (który już wcześniej został obliczony).
	\end{itemize}
	\item Jeżeli obliczony zysk jest większy niż dotychczasowa maksymalna wartość $q$, następuje aktualizacja:
	\begin{itemize}
		\item $q$ przyjmuje nową, większą wartość zysku,
		\item $s[j]$ zapisuje długość cięcia $i$ jako optymalne pierwsze cięcie dla pręta o długości $j$.
	\end{itemize}
\end{enumerate}

\large{Pętla kończy się, gdy wszystkie możliwe długości cięć zostaną przeanalizowane, a zmienna $q$ zawiera maksymalny możliwy zysk dla pręta o długości $j$.}

\section{LCS - problem znajdowania najdłuższego wspólnego podciągu}

\large{Przy analizie algorytmu rozwiązującego problem znajdowania najdłuższego wspólnego podciągu przyjrzymy się implementacji rekurencyjnej ze spamiętywaniem oraz wersji iteracyjnej.}

\subsection{rekurencyjna ze spamiętywaniem}

\begin{lstlisting}[language=C++, caption={Kod funkcji znajdowania nadłuższego wspólnego podciągu w wersji rekurencyjnej}, numbers=left, frame=single, xleftmargin=20pt]
int LCS_R(const string &w1, const string &w2, 
   int m, int n, int** s) {
	if (s[m][n] != -1) {
		return s[m][n];
	}
	if (m == 0 || n == 0) {
		return s[m][n] = 0;
	}
	if (w1[m - 1] == w2[n - 1]) {
		return s[m][n] = 1 + LCS_R(w1,
		 w2, m - 1, n - 1, s);
	}
	return s[m][n] = max(LCS_R(w1, w2, m,
	 n - 1,s), LCS_R(w1, w2, m - 1, n,s));
}
\end{lstlisting}

\large{Argumentami funkcji w implementacji rekurencyjnej są jakieś litery oraz ich długości $w1$, $w2$ i odpowiednio $m$ oraz $n$, a także tablica $s[]$ mająca za zadanie przechowywać maksymalne długości podciągów. Jak w każdej rekurencyjnej funkcji mamy warunki początkowe. Najważniejszą częścią tej implementacji jest podzielenie sytuacji na dwa przypadki - w której mamy styczność z takimi samymi ostatnimi znakami wyrazu (dodajemy $1$ i wywołujemy rekurencyjne dla pozstałych) lub są różne. W takim przypadku obliczamy dwie możliwości (z pominięciem ostatniego znaku $w1$ lub $w2$) i zwracamy wartość maksymalną. Tablica $s[]$ jest zwracana i to ona odpowiada za funkcję spamiętywania.}

\subsection{iteracyjna}
\begin{lstlisting}[language=C++, caption={Kod funkcji znajdowania nadłuższego wspólnego podciągu w wersji iteracyjnej}, numbers=left, frame=single, xleftmargin=20pt]
char** LCS_I(const string &w1, const 
   string &w2, int& dlugosc) {
	int m = w1.length();
	int n = w2.length();
	
	int** c = new int*[m + 1];
	char** b = new char*[m + 1];
	for (int i = 0; i <= m; i++) {
		c[i] = new int[n + 1];
		b[i] = new char[n + 1];
	}
	
	for (int i = 0; i <= m; i++) {
		c[i][0] = 0;
	}
	for (int j = 0; j <= n; j++) {
		c[0][j] = 0;
	}
\end{lstlisting}

\large{Przypadek iteracyjny został napisany na podstawie pseudokodu, więc nie jest on zbytnio ciekawym kodem. W mojej implementacji mamy tylko trzy argumenty: $w1$, $w2$ oraz zmienna $dlugosc$. Dopiero wewnątrz funkcji definiuje zmienne $m$ oraz $n$ jak również tworzę dwie pomocnicze tablice: $b[]$ odpowiedzialna za przechowywanie kierunków do odtworzenia LCS oraz $c[]$ przechowująca długość ciągu.}

\begin{lstlisting}[language=C++, caption={Kod funkcji znajdowania nadłuższego wspólnego podciągu w wersji iteracyjnej}, numbers=left, frame=single, xleftmargin=20pt]
if (w1[i - 1] == w2[j - 1]) {
	c[i][j] = c[i - 1][j - 1] + 1;
	b[i][j] = '\\';
} else if (c[i - 1][j] >= c[i][j - 1]) {
	c[i][j] = c[i - 1][j];
	b[i][j] = '|';
} else {
	c[i][j] = c[i][j - 1];
	b[i][j] = '-';
	}
}
dlugosc = c[m][n];
\end{lstlisting}

\large{przeanalizujmy najciekawszy część tego algorytmu:}

\begin{itemize}
	\item Dla każdej pary $(i, j)$, gdzie $1 \leq i \leq m$ oraz $1 \leq j \leq n$:
	\begin{enumerate}
		\item Jeśli znaki w obu ciągach są równe, tj. $w1[i-1] = w2[j-1]$, to:
		\[
		c[i][j] = c[i-1][j-1] + 1
		\]
		oraz ustawiamy kierunek na ukośny
		\item Jeśli znaki są różne, porównujemy wartości:
		\[
		c[i-1][j] \quad \text{oraz} \quad c[i][j-1]
		\]
		\begin{itemize}
			\item Jeśli $c[i-1][j] \geq c[i][j-1]$, to:
			\[
			c[i][j] = c[i-1][j]
			\]
			oraz ustawiamy kierunek pionowy
			\item W przeciwnym przypadku:
			\[
			c[i][j] = c[i][j-1]
			\]
			oraz ustawiamy kierunek poziomy
		\end{itemize}
	\end{enumerate}
\end{itemize}

\large{Na sam koniec przypisujemy do argumentu $dlugosc$ wartość najdłuższego wspólnego podciągu. Funkcja zwraca $b[]$.}

\end{document}